\documentclass[11pt,twoside,openright]{book}

\usepackage[T1]{fontenc}
\usepackage[utf8]{inputenc}
\usepackage{amsmath,amssymb,amsthm,mathtools,bm}
\usepackage{mathpazo}
\usepackage{microtype}
\usepackage[inner=1.35in,outer=1.1in,top=1.25in,bottom=1.35in]{geometry}
\usepackage{booktabs}
\usepackage{fancyhdr}
\usepackage[hidelinks]{hyperref}
\usepackage{cleveref}

\setlength{\parindent}{1.4em}
\setlength{\parskip}{0pt}
\linespread{1.07}
\emergencystretch=3em

\pagestyle{fancy}
\fancyhf{}
\fancyhead[LE]{\small\thepage\quad\textsc{A Cinema of Individual Editions}}
\fancyhead[RO]{\small\nouppercase{\leftmark}\quad\thepage}
\renewcommand{\headrulewidth}{0.3pt}
\fancypagestyle{plain}{\fancyhf{}\fancyfoot[C]{\small\thepage}\renewcommand{\headrulewidth}{0pt}}

\theoremstyle{plain}
\newtheorem{theorem}{Theorem}[chapter]
\newtheorem{proposition}[theorem]{Proposition}
\newtheorem{lemma}[theorem]{Lemma}
\newtheorem{corollary}[theorem]{Corollary}
\theoremstyle{definition}
\newtheorem{definition}[theorem]{Definition}
\newtheorem{example}[theorem]{Example}
\theoremstyle{remark}
\newtheorem{remark}[theorem]{Remark}

\newcommand{\Capac}{\mathrm{Cap}}
\newcommand{\Rob}{\mathrm{Rob}}
\newcommand{\Inv}{\mathrm{Inv}}
\DeclareMathOperator{\loss}{loss}

\title{\Huge\bfseries A Cinema of Individual Editions\\[1.2ex]
\large Chapter Two: Perceptual Equivalence and the Measurement of Difference\\[0.4ex]
\large Working Draft}
\author{Flyxion\\[0.3ex]\small Independent Researcher}
\date{November 2026}

\begin{document}

\frontmatter
\maketitle
\tableofcontents
\mainmatter
\setcounter{chapter}{1}

\chapter{Perceptual Equivalence and the Measurement of Difference}
\label{ch:perceptual}

Chapter~1 counted editions by assuming that perceptual classes exist and that two editions are equivalent exactly when they share a class at every slot. The assumption made the counts exact, and every later chapter has used it. This chapter replaces it with an account of dissimilarity and says what the replacement permits. The account separates three things that a viewer can notice, the presence of a difference, its bearing on what the story means, and its effect on how the film feels, and gives each its own pseudometric. It shows what happens to the counts when the equivalence is replaced by a threshold, where the replacement fails to behave like an equivalence, and how the packing and covering numbers that take its place depend on the question the production wants to answer. It then defines the comparison across productions that the worldwide promise of Chapter~11 will need.

The chapter proves what the formalism permits and marks the places where evidence about viewers has to take over from mathematics. The lawn scene supplies the numbers, and the choices of pseudometric used for them are illustrations. None of the figures below is a measurement of any audience.

\section{What Exact Equivalence Assumed}\label{sec:assumed}

Exact equivalence says that a viewer cannot tell two editions apart if and only if they have the same class at every slot. It has two halves. If the classes agree at every slot, the viewer cannot tell the editions apart, which says that the class partition is fine enough to hide nothing a viewer could notice. If some class differs, the viewer can tell them apart, which says that every class boundary is perceptible. Both halves are strong, and the book needs only that the counts be computed from some partition. What the first half guarantees is that the editions counted as distinct are not secretly the same, and what the second half guarantees is that no distinct editions are being merged. A production that wanted to promise that no two viewers get the same film would rely on the second half, and a production that wanted to promise each viewer something perceptibly different would rely on the first.

A more careful account does not partition editions into classes in the first place. It gives each pair of editions a number measuring how different they are.

\section{Three Notions of Difference}\label{sec:three2}

A \emph{pseudometric} on editions is a function $d$ with $d(x,x)=0$, symmetry, and the triangle inequality, without requiring that $d(x,y)=0$ implies $x=y$. The editions at distance zero from one another form an equivalence class, and $d$ is a metric on the classes. A pseudometric therefore carries exact equivalence inside it: the equivalence is the relation $d=0$.

Three pseudometrics are distinguished by what they measure.

\begin{definition}[Detection, comprehension, affect]\label{def:three}
The \emph{detection} dissimilarity $d_{\mathrm{det}}$ measures whether a viewer can notice, on comparison, that two editions differ. The \emph{comprehension} dissimilarity $d_{\mathrm{comp}}$ measures whether the two editions support different inferences about what happened in the story. The \emph{affective} dissimilarity $d_{\mathrm{aff}}$ measures whether they produce different responses of feeling, tension, or attention.
\end{definition}

The three can disagree. Two editions that differ only in the framing of a shot are detectably different and have the same comprehension. Two that differ in whether grass reappears have different comprehension, and may or may not feel different. Two that differ in the direction of a long take differ in affect and may not differ in what they say happened.

For the lawn scene the following instances are used, as illustrations. Detection is the Hamming dissimilarity on class tuples, the number of slots at which two editions use different classes. Affect is the Hamming dissimilarity on direction sequences, so that two classes of the same screen direction are affectively alike. Comprehension is the Hamming dissimilarity on the pattern of stale classes, because the only narrative information that the scene carries is the state of the grass, and an edition conveys a different story only where a stale class misstates it. The pattern of stale slots records where the film disagrees with the script, which is not necessarily what story it conveys to a viewer, and the choice is an illustration for that reason. The editions at distance zero under each of these are then counted. \Cref{tab:notions} gives the numbers of editions that remain distinguishable.

\begin{table}[ht]
\centering
\begin{tabular}{@{}llrrr@{}}
\toprule
Production & Rules & Detection & Affect & Comprehension\\
\midrule
A & direction rule & 494 & 50 & 8\\
A & plus agreement rule & 238 & 40 & 1\\
B & direction rule & 320 & 2 & 5\\
B & plus agreement rule & 96 & 2 & 1\\
C & direction rule & 2 & 1 & 2\\
C & plus agreement rule & 1 & 1 & 1\\
\bottomrule
\end{tabular}
\caption{Number of classes of editions at distance zero under each of three illustrative pseudometrics, for the lawn scene.}\label{tab:notions}
\end{table}

The table carries a result that is worth stating plainly. Once the agreement rule is in force, every admissible edition of every production tells the same story, so the lawn scene has exactly one edition in the sense of comprehension. All of the variation that the book counts is variation of detection and, to a lesser extent, of affect. This is not a defect of the counting. It is what it means for the rules to be correct: the agreement rule is the requirement that editions not differ in what they say. The variation that remains is variation in how it is said. Production B has two affective classes, all its slots committed to the left or all to the right, and $96$ detectable ones.

\section{What an Experiment Would Have to Measure}\label{sec:experiment}

Each of the three pseudometrics names a quantity that a study could in principle estimate, and none has been estimated here.

For detection, a study would present pairs of editions to viewers and record how often the viewer reports a difference, with the pairs chosen to span the slots and classes of the production. The result is a measure of how detectable each class pair is, and the question for the formalism is whether the resulting dissimilarity is close to a sum over slots, as the Hamming choice assumes. A viewer might detect one large difference among five small ones that it ignores, and then the dissimilarity is not additive. Raw pairwise detection frequencies also need not satisfy the triangle inequality. If viewer data resist a pseudometric model, the formalism can be kept by taking the shortest-path closure of the measured dissimilarities, which is the largest pseudometric below them, or by treating the pairs that viewers can tell apart as a graph and asking for independent sets in it, which loses the metric structure and keeps the packing question.

For comprehension, a study would ask viewers questions about the story after watching one edition, and compare the answers across editions. The measure is whether the distribution of answers differs, and the interesting cases are those in which two editions give different answers to a question that the script history answers one way. These are the cases that the agreement rule exists to prevent.

For affect, a study would collect continuous or retrospective ratings of tension, interest, or emotion during viewing, and compare them across editions. This is the hardest of the three, since viewers differ in how a given shot makes them feel and the response to an edition depends on what the viewer expects of the film.

For each, three numbers have to be fixed: a tolerance below which the difference is treated as none, the population over which the measure is taken, and the stability of the measure across viewers. The formalism of this chapter depends on the tolerance, and the later chapters on the population.

\section{Thresholds Are Not Equivalences}\label{sec:thresholds}

If the tolerance is $\varepsilon$, two editions can be called \emph{$\varepsilon$-similar} when their dissimilarity is at most $\varepsilon$. For $\varepsilon=0$ this is the exact equivalence. For $\varepsilon>0$ it is reflexive and symmetric and it is not transitive.

\begin{proposition}[Failure of transitivity]\label{prop:transitivity}
Let $d$ be a pseudometric and $\varepsilon>0$. The relation $d\le\varepsilon$ is reflexive and symmetric and need not be transitive: an edition $x$ can be $\varepsilon$-similar to $y$, and $y$ to $z$, while $x$ and $z$ are not. The triangle inequality gives only $d(x,z)\le2\varepsilon$.
\end{proposition}

\begin{proof}
The triangle inequality bounds $d(x,z)$ by $d(x,y)+d(y,z)\le2\varepsilon$ and no better in general. A line with points at $0$, $\varepsilon$ and $2\varepsilon$ has $0\sim\varepsilon$ and $\varepsilon\sim2\varepsilon$ and $0\not\sim2\varepsilon$.
\end{proof}

The lawn scene shows how much of the failure there is. Take the $238$ editions of production A after the agreement rule, under Hamming dissimilarity on class tuples. At $\varepsilon=1$, two editions are similar if they differ at most at one slot. There are $10{,}040$ triples in which one edition is similar to two others, and in $8{,}228$ of them the two others are not similar to each other. At $\varepsilon=2$ the numbers are $263{,}740$ and $147{,}328$. In both cases the similarity graph is connected, so the transitive closure of similarity is a single class. A production that tried to turn the threshold into an equivalence by closing it would conclude that all $238$ editions are one film. For production B, with $96$ editions, the numbers at $\varepsilon=1$ are $1{,}440$ and $1{,}344$, and at $\varepsilon=2$ they are $18{,}240$ and $8{,}352$, and the closure has two components, the all-left editions and the all-right ones.

The failure of transitivity means that the count of editions cannot be generalized by counting classes of a similarity relation. Two replacements remain, and they ask different questions.

\section{Packing and Covering}\label{sec:packcover}

The \emph{packing number} $M_\varepsilon$ is the size of a largest set of editions that are pairwise more than $\varepsilon$ apart, as in Chapter~1. The \emph{covering number} $N_\varepsilon$ is the size of a smallest set of editions such that every edition is within $\varepsilon$ of one of them.

The two answer different questions. The packing number is the number of viewers who can each be given an edition that is perceptibly different from every other viewer's, which is the sustainable audience of a promise of distinctness. The covering number is the number of editions that suffice to represent every possible edition within the tolerance, which is the number of distinct films the audience could be shown if the production were content to show each viewer something within $\varepsilon$ of the film the viewer's assignment represents. It is the quantity that a compression of the space of editions at tolerance $\varepsilon$ needs, and the logarithm $\log_2N_\varepsilon$ is the number of bits that an index into such a compression takes.

\begin{proposition}[Packing and covering]\label{prop:pc}
For every pseudometric on a finite set and every integer tolerance $\varepsilon$ of a Hamming-type dissimilarity, $N_\varepsilon\le M_\varepsilon\le N_{\lfloor\varepsilon/2\rfloor}$. For a general pseudometric the last term is $N_{\varepsilon/2}$.
\end{proposition}

\begin{proof}
A largest set of editions that are pairwise more than $\varepsilon$ apart is maximal, so every other edition is within $\varepsilon$ of it, and the set covers at radius $\varepsilon$, which gives $N_\varepsilon\le M_\varepsilon$. For the second inequality, suppose a set of editions is pairwise more than $\varepsilon$ apart, and cover the space by balls of radius $\lfloor\varepsilon/2\rfloor$. Two members of the set in one ball would be within twice that radius of each other, at most $\varepsilon$, a contradiction. So each ball holds at most one member of the set, and there are at least as many balls as members.
\end{proof}

\Cref{tab:packcover} gives the numbers for the lawn editions after the agreement rule, under Hamming dissimilarity on class tuples. The packing numbers were computed exactly by integer programming. The covering numbers were computed by integer programming within a time limit. They are exact where the bounds meet, and for production A at $\varepsilon=1$ the solver proved only that the number lies between $25$ and $31$.

\begin{table}[ht]
\centering
\begin{tabular}{@{}llrrrr@{}}
\toprule
Production & $\varepsilon$ & Packing $M_\varepsilon$ & Covering $N_\varepsilon$ & $\log_2 M_\varepsilon$ & $\log_2 N_\varepsilon$\\
\midrule
A & 0 & 238 & 238 & 7.9 & 7.9\\
A & 1 & 54 & 25 to 31 & 5.8 & 4.6 to 5.0\\
A & 2 & 16 & 8 & 4.0 & 3.0\\
B & 0 & 96 & 96 & 6.6 & 6.6\\
B & 1 & 32 & 16 & 5.0 & 4.0\\
B & 2 & 12 & 6 & 3.6 & 2.6\\
\bottomrule
\end{tabular}
\caption{Packing and covering numbers of the admissible editions after the agreement rule, under Hamming dissimilarity on class tuples.}\label{tab:packcover}
\end{table}

The numbers show how sharply the audience that a promise can sustain depends on the tolerance. Production A can serve $238$ viewers with exactly distinct editions and $54$ with editions that differ at two or more slots, and $16$ with editions that differ at three or more. The covering numbers say how many reference editions would stand in for the whole space: eight at $\varepsilon=2$ for A. This is the same arithmetic that Chapter~1 used for its rate, with a different numerator. The rate of that chapter counts exactly distinguishable editions. The covering count is the quantity for a production that is content with distinguishability up to a tolerance.

The two numbers are not interchangeable, and the choice between them is part of what a production promises. A promise of pairwise distinctness is a packing promise and its audience is bounded by $M_\varepsilon$. A promise that every viewer's edition is within $\varepsilon$ of an agreed reference is a covering promise and costs $N_\varepsilon$ references. A production that confuses them will state a sustainable audience of $N_\varepsilon$ when the audience it can sustain is $M_\varepsilon$, and in \Cref{tab:packcover} the two differ by a factor of about two.

\section{Across Productions}\label{sec:across}

The worldwide promise of Chapter~1 compares editions from different productions, and the pseudometric has to be extended to editions that were not assembled from the same pool. Two cases arise.

When the productions share no footage, so that no class at any slot of one is a class at the same slot of the other, every pair of editions from different productions differs at every slot. Under Hamming dissimilarity the cross distance is the number of slots, here $5$, and the comparison is trivial for every tolerance below $5$.

\begin{proposition}[Disjoint footage composes]\label{prop:disjoint}
If two productions share no class at any slot, then for every $\varepsilon<m$ the packing number of the union of their editions is the sum of the packing numbers, and exact inequivalence composes: the number of exactly distinct editions of the union is the sum of the two.
\end{proposition}

\begin{proof}
Every cross pair is at distance $m>\varepsilon$, so a separated set in each production remains separated when combined, and any separated set of the union splits into separated sets of the two. The same holds with $\varepsilon=0$.
\end{proof}

When productions share footage the cross distance can be small, and zero. To see the effect, suppose that classes with the same direction, status, and index within their group at a slot are the same footage in productions A and B. Then $4$ of the class tuples of the two productions after the agreement rule are identical. The exact count of the union is $330$, and not the sum $334$, and the effect appears at every tolerance. At $\varepsilon=1$ the packing numbers are $54$ and $32$ for A and B and the packing number of the union is $85$ and not $86$. At $\varepsilon=2$ the numbers are $16$, $12$ and $27$ against $28$. A pair of viewers in two locales who receive editions that share all their footage cannot be given editions that differ.

The calculation is small, but its form is general. Shared footage, such as an establishing shot used in every locale, makes the global audience smaller than the sum of the local audiences, and a production that shares more footage saves storage and loses distinctness. This is the rate trade-off of Chapter~1 at the scale of a release.

\section{What the Formalism Permits}\label{sec:permits}

The formalism permits the following. Given a pseudometric, the equivalence it defines is exact and the counts of Chapter~1 follow. Given a tolerance, the packing number bounds the audience of a distinctness promise and the covering number gives the size of a compression, with the inequality of \Cref{prop:pc} relating them. Given two productions, the combined audience is the sum when they share no footage and is smaller when they do. All of these hold for any pseudometric, and none of them needs a theory of what viewers perceive.

The formalism does not say which pseudometric describes a given audience, what tolerance is perceptible, or whether the three notions of difference can be treated separately. It does not say whether detection is additive over slots, or whether viewers respond to the same differences the pseudometric counts. These are questions about people, and the answers have to come from studies that have not been done. Until they are done, the numbers in this chapter show how the quantities behave and not what they are. The remaining chapters take the Hamming dissimilarity on class tuples as a stand-in whenever a tolerance is needed.

The scripts \texttt{lawnmodel.py} and \texttt{verify\_ch02.py} reproduce each figure of this chapter.

\begin{thebibliography}{99}

\bibitem{flyxion2026gf}
Flyxion.
\newblock \emph{Generative Film: Categories, Sheaves, and Types for Cinema}.
\newblock Unpublished monograph, October 2026.

\end{thebibliography}

\end{document}
